锂盐是雅化利润弹性的主要来源，也是估值中最容易产生伪精度的部分。公司在 H1 预告中确认了价格、销量和销售均价同步改善，但没有披露 H1 的实际销量、实现价格、单位成本或分部毛利。因此，我们用 2025 年经审计经营基线建立情景，而不把现货价格或设计产能直接换算成 2026 年 EPS。

\section{2025 年的经营底座：已知的是产品篮子，不是 2026 年单位利润}

2025 年锂产品收入为 48.24 亿元、销量为 6.95 万吨，折算的综合收入代理为 6.94 万元/吨；成本为 42.34 亿元，毛利率为 12.22\%。这组数值可以定义周期底部后的历史产品篮子，却不能证明 2026 年每吨实现价格、矿石成本或库存收益。尤其是公司尚未披露实际自供比例、KMC/李家沟的实际交付和现金成本，基准情景不因此给予成本优势溢价。

\begin{exhibitbox}[表 3-1：锂盐历史经营锚——模型使用的已审计起点]
\begin{tabularx}{\exhibitboxwidth}{L{4.2cm}R{2.8cm}X}
\toprule
指标（2025A） & 数值 & 模型用途与边界 \\
\midrule
锂产品收入 & 48.24 亿元 & 建立产品篮子收入基数，不等于 2026 年收入 \\
锂产品销量 & 69,493.81 吨 & 历史销售量锚；不以近 13 万吨设计产能替代 \\
综合收入代理 & 6.94 万元/吨 & 收入除销量；非公司披露的单一产品 ASP \\
锂产品成本 & 42.34 亿元 & 历史成本锚；非 2026 年单位现金成本 \\
毛利率 / 毛利代理 & 12.22\% / 0.85 万元/吨 & 识别周期弹性，不直接外推至全年 \\
\bottomrule
\end{tabularx}
\end{exhibitbox}
\sourcenote{公司《2025 年年度报告》分产品收入、成本及销售量披露；综合收入与毛利代理为 AStock 用上述审计数复算。}

这一边界直接影响估值：如果 2026 年的改善主要是价格、低价库存或一次性的销售节奏，较高的 H1 利润并不应获得长期 PE；如果销量、实现价格、成本和现金同时改善，才说明周期反弹正在形成可持续的单位经济。

\section{自有情景：H1 预告锚定上半年，H2 不作 Q2 年化}

我们的 2026 年模型以 H1 归母 11.0/12.0/13.0 亿元作为熊/基准/牛的上半年锚，再分别假设 H2 归母 4.0/8.1/11.7 亿元。基准情景的 H2 利润低于 Q2 预告中位数，反映我们没有把单季高盈利机械延长到全年。该处理使模型能够回答“当前价格需要什么持续性”，而不把尚未公布的 H1 分部数据伪装成事实。

\begin{exhibitbox}[表 3-2：锂盐情景与合并盈利桥]
\begin{tabularx}{\exhibitboxwidth}{L{3.0cm}R{2.2cm}R{2.2cm}R{2.2cm}X}
\toprule
2026E 自有情景 & 熊 & 基准 & 牛 & 估值含义 \\
\midrule
锂盐销量 & 7.5 万吨 & 9.0 万吨 & 10.0 万吨 & 仅为情景输入，非公司指引 \\
实现收入代理 & 9.5 万元/吨 & 10.5 万元/吨 & 11.2 万元/吨 & 以产品篮子代理，非单一产品 ASP \\
锂产品毛利率 & 25\% & 32\% & 35\% & 未给自供矿、套保或库存额外信用 \\
H1 归母锚 & 11.0 亿元 & 12.0 亿元 & 13.0 亿元 & 公司业绩预告区间内的情景锚 \\
H2 归母假设 & 4.0 亿元 & 8.1 亿元 & 11.7 亿元 & 不以 Q2 单季利润年化 \\
2026E 归母 / EPS & 15.0 亿元 / 1.30 元 & 20.1 亿元 / 1.74 元 & 24.7 亿元 / 2.14 元 & 供第八章估值使用 \\
\bottomrule
\end{tabularx}
\end{exhibitbox}
\sourcenote{公司《2026 年半年度业绩预告》；AStock 情景模型。情景为公开资料下的分析假设，正式 H1 披露后必须重估。}

基准情景的核心不是“锂价继续上涨”，而是销量、实现价格和毛利率共同达到表 3-2 的中位假设。H1 正式报告若不能给出足够的收入、库存和现金线索来支持这一组合，模型应向熊情景迁移；反之，若单位经济和现金同步兑现，牛情景才有资格获得更高概率。

\section{资源与客户：保留供应保障框架，不预支成本优势}

公司披露持有能投锂业 37.25\% 股权并享有李家沟锂精矿优先供应权，也披露 KMC、李家沟及多个锂盐基地的资源/产能布局。它们说明公司存在供应保障路径，却没有在截止日披露雅化 2026 年实际取得的数量、价格、自供比例、回收率或单位成本。具名客户同样证明资质与渠道，而不证明客户份额、剩余订单或 2026 年收入。

\begin{sourcequalitybox}[估值处理]
李家沟/KMC 的实际供给和成本、客户的订单与份额、锂盐基地利用率、实际 ASP、单位成本、套保和硫化锂商业化均为零正向估值信用。它们只构成牛情景必须额外验证的上行来源，而不是基准 EPS 的隐含前提。
\end{sourcequalitybox}
